% ======================================================================
% Sections 6–8 — Discussion, Limitations, Conclusion
% Revised from submitted analysis/conclusion + final Results framing
% ======================================================================
\section{Discussion}\label{sec:discussion}

\subsection{Why the Physical Graph Can Be Harmful}\label{sec:disc-dense}

The physical road-network graph is dense relative to the number of sensors
(on the order of $10^3$ off-diagonal entries for $207$ Los-loop detectors).
Under repeated neighborhood aggregation this density can oversmooth node
representations, so that a graph-based forecaster is worse than the same
backbone with identity adjacency. The consistent gap between Physical and
NoSpatial on both benchmarks and both backbones is consistent with that
reading. It also changes how GSL should be judged: improving on the
physical graph is not sufficient evidence that learned structure is useful,
because discarding the graph already helps.

Learned contemporaneous graphs occupy an intermediate position --- less
harmful than the physical graph, but not better than using no graph. What
they appear to provide is a less harmful neighborhood, not enough
additional signal to beat identity under these protocols.

\subsection{When Learned Structure Helps}\label{sec:disc-help}

The positive Los-loop result is specific: multi-lag DAGMA graphs consumed
in a lag-aligned way inside T-GCN. Three observations support a
structure-learning interpretation rather than a pure sparsity
interpretation. First, at a matched $30$-edge budget, random and
correlation-placed graphs are worse than no graph, while DAGMA-placed lag
graphs are better (Los-loop PH1 only). Second, the multi-lag construction
recovers substantially more lagged structure on Los-loop than on SZ-Taxi
($30$ vs.\ $2$ slots). Third, the same multi-lag artifacts fail when merged
into a static union for GCN, so estimation alone is not enough; the
backbone must use the structure in a way that matches how it was organized
(lag blocks).

The GCN/T-GCN contrast for contemporaneous graphs (cGSL helps GCN more than
T-GCN) is an aggregation-compatibility observation of this study. We do not
interpret the contemporaneous DAG as a temporal graph
(Section~\ref{sec:gsl_contemp}); temporal organization of learned structure
is carried by the multi-lag construction only.

\subsection{What the Learned Graph Represents}\label{sec:disc-interpretation}

We treat the learned adjacencies as \emph{statistical dependency}
estimates under a linear DAGMA fit. Acyclicity is an optimization
constraint, not evidence of causal traffic mechanisms. The lag-structured
reading of the multi-lag blocks is consistent with the input construction
$Z=[x(t-3),\ldots,x(t)]$ and with the consumption experiments, but it is
not independently validated against ground-truth propagation. The gated
variant changes how static lag graphs are used per node and timestep; it
does not produce a time-varying edge set.

\subsection{Dataset Dependence}\label{sec:disc-dataset}

SZ-Taxi and Los-loop differ in sampling interval, sensor count, and road
environment. The $15$-minute Los-loop control (Section~\ref{sec:res-boundaries})
shows that coarser sampling alone does not suppress the multi-lag gain.
The more concrete difference in these experiments is how much structure the
estimator recovers: a nearly empty multi-lag fit on SZ-Taxi is associated
with results close to the graph-free baseline. This is a scope condition on
GSL for traffic, not a failure of the general idea that adjacency can be
estimated from data.

\section{Limitations}\label{sec:limitations}

\begin{itemize}
\item \textbf{Statistical power.} Main results use five training seeds.
Sample standard deviations and paired win counts are reported; with $n=5$
the exact Wilcoxon test cannot fall below $p=0.0625$, so we avoid definitive
significance claims.
\item \textbf{Control scope.} Matched-sparsity and capacity controls exist
only for Los-loop at PH1. There is no sparsified physical graph, no
$\lambda$/threshold sweep, and no SZ-Taxi sparsity control.
\item \textbf{Datasets and horizons.} Two benchmarks; $\mathrm{PH}\le 4$ at
native sampling. The $15$-minute Los-loop study is a resolution control, not
evidence for PH5--8 at $5$-minute sampling.
\item \textbf{Estimator assumptions.} Linear DAGMA with $\ell_2$ score and
acyclicity regularizer; nonlinear dependency and misspecified lag depth
$L=3$ are not examined. Multi-lag fitting is expensive for large $N$.
\item \textbf{Static graphs.} Learned adjacencies are fitted once. Mix
varies \emph{use} of those graphs, not the edge set; time-varying structure
learning remains future work.
\item \textbf{Backbones.} Evaluation is limited to GCN and T-GCN. Larger
attention or transformer traffic models are outside the present study.
\end{itemize}

\section{Conclusion and Future Works}\label{sec:conclusion}

This paper revisits graph structure learning for traffic forecasting under
a protocol that includes physical, graph-free, contemporaneous learned, and
multi-lag learned adjacencies on two public benchmarks. The main findings
are:

\begin{enumerate}
\item A dense physical road-network graph can be a poor default for GCN and
T-GCN at short horizons; a graph-free identity control is a necessary
reference when evaluating learned structure.
\item A single contemporaneous DAGMA graph improves on the physical graph in
many settings but does not beat the graph-free control on these datasets.
\item Multi-lag learned graphs help on Los-loop when consumed in a
lag-aligned way inside T-GCN; per-node gating (T-GCN-MultiGSL-Mix) is the
strongest configuration ($14.5\%$ relative RMSE reduction vs.\ the
graph-free baseline at PH1; up to $43.0\%$ vs.\ the physical baseline
across horizons). Matched-budget controls show that sparsity alone does not
explain this gain on the studied cell.
\item The benefit is dataset-dependent. On SZ-Taxi, where the multi-lag fit
is nearly empty, learned variants remain close to the graph-free baseline.
A $15$-minute resampling of Los-loop retains a large relative gain, so
coarser sampling alone does not explain the SZ-Taxi result.
\item Learned graphs are interpreted as statistical dependency structure,
not causal maps of traffic influence.
\end{enumerate}

Future work includes longer horizons at $5$-minute resolution, additional
networks, sparsified-physical controls and $\lambda$/threshold sensitivity,
other backbones, hybrid combinations with attention
\citep{Sun2023Attention}, and genuinely time-varying graph estimation
(e.g.\ incremental or sliding-window structure learning) rather than gated
use of a fixed graph. Overall, learning the adjacency from traffic data can
improve forecasting when the recovered structure is informative and is used
in a manner consistent with how it was estimated; it is not a universal
replacement for either the physical graph or a graph-free model.
